\documentclass[a11paper]{article}

\usepackage{karnaugh-map}
\usepackage{tabularx}
\usepackage{libs/titlepage}
\usepackage{amsmath}
\usepackage{amsthm}
\usepackage{libs/document}
\usepackage{booktabs}
\usepackage{multicol}
\usepackage{float}

\usepackage{witharrows}
\usepackage{enumitem}

\usepackage[usenames,dvipsnames]{xcolor}

% \renewcommand{\not}[1]{\overline{#1}}
\renewcommand{\not}[1]{#1^{\prime}}
\newcommand{\rhs}[1]{#1^{\prime}}
\newcommand{\xor}{\oplus}
\newcommand{\Cuz}[1]{\Arrow{\footnotesize #1}}

% \item $A+B=B+A \hspace{\fill}\text{commutativity}$
\usepackage{enumitem}       % customizable list environments


\title{Devoir}

\class{Logique Combinatoire}
\classnb{APP1}

\teacher{Marwan Besrour \& Gabriel Bélanger}

\author{Benjamin Chausse --- CHAB1704}

\begin{document}
\maketitle

\section{Démontrer la solution la plus simple au segment G d'un afficheur 7 segment}

Commençons par dresser la table de vérité du segment G de l'afficheur. À partir
de celle-ci, il sera possible de trouver la solution en créant une table de
Karnaugh à partir de cette table.

\begin{table}[H]
  \centering
  \footnotesize
  \caption{Table de vérité du segment G}
  \begin{tabular}{lllllc}
    \toprule
    Dec. & $D_3$ & $D_2$ & $D_1$ & $D_0$ & \textbf{G} \\
    \midrule
    0  & 0 & 0 & 0 & 0 & 0 \\
    1  & 0 & 0 & 0 & 1 & 0 \\
    2  & 0 & 0 & 1 & 0 & 1 \\
    3  & 0 & 0 & 1 & 1 & 1 \\
    4  & 0 & 1 & 0 & 0 & 1 \\
    5  & 0 & 1 & 0 & 1 & 1 \\
    6  & 0 & 1 & 1 & 0 & 1 \\
    7  & 0 & 1 & 1 & 1 & 0 \\
    8  & 1 & 0 & 0 & 0 & 1 \\
    9  & 1 & 0 & 0 & 1 & 1 \\
    10 & 1 & 0 & 1 & 0 & X \\
    11 & 1 & 0 & 1 & 1 & X \\
    12 & 1 & 1 & 0 & 0 & X \\
    13 & 1 & 1 & 0 & 1 & X \\
    14 & 1 & 1 & 1 & 0 & X \\
    15 & 1 & 1 & 1 & 1 & X \\
    \bottomrule
  \end{tabular}
\end{table}

\begin{figure}[H]
  \centering
    \begin{karnaugh-map}[4][4][1][$D_0$][$D_1$][$D_2$][$D_3$]
        \manualterms{
          0,0,1,1,
          1,1,1,0,
          1,1,X,X,
          X,X,X,X}
        \implicant{2}{10}
        \implicant{4}{13}
        \implicant{12}{10}
        \implicantedge{3}{2}{11}{10}
    \end{karnaugh-map}
  \caption{Table de Karnaugh du segment G}
\end{figure}

En représentant de façon formelle ce que nous révèle la table de Karnaugh
nous obtenons l'équation que nous avions à démontrer:

\begin{equation}
  G = \underbrace{D_3}_{\textit{\textcolor{YellowOrange}{Jaune}}}+
      \underbrace{D_2\not{D_1}}_{\textit{\textcolor{ForestGreen}{Vert}}}+
      \underbrace{\not{D_2}D_1}_{\textit{\textcolor{RoyalBlue}{Bleu}}}+
      \underbrace{D_1\not{D_0}}_{\textit{\textcolor{Mahogany}{Rouge}}}
\end{equation}

\section{Prouver que l'égalité qui suit est vraie par l'algèbre de Boole}

Équation initiale:

\begin{equation}
  \tiny
  \not{D_3} \not{D_2} D_1 +
  \not{D_3} D_1 \not{D_0} +
  D_3 \not{D_2} \not{D_1} +
  \not{D_3} D_2 \not{D_1} =
  D_3 \left[ (D_2 \xor D_1) + \not{( \not{D_1} + D_0)} \right] +
  D_3(\not{D_1}\not{(D_2\not{D_1})}) +
  (D_1\xor D_3)(D_1\xor D_2)(D_1(D_2\not{D_1} + \not{D_2}\not{D_1}))
\end{equation}

Afin de simplifier la résolution (et alléger la lecture), les termes de
l'équation sont substitués par de plus petits segments représentés par des
variables:

\begin{align}
  A &= \not{D_3} \left[ (D_2 \xor D_1) + \not{( \not{D_1} + D_0)} \right] \\
  B &= D_3(\not{D_1}\not{(D_2\not{D_1})}) \\
  C &= (D_1\xor D_3)(D_1\xor D_2)(D_1(D_2\not{D_1} + \not{D_2}\not{D_1})) \\
  \not{D_3} \not{D_2} D_1 + \not{D_3} D_1 \not{D_0} + D_3 \not{D_2} \not{D_1} + \not{D_3} D_2 \not{D_1} = R &= A+B+C
\end{align}

\subsection{Simplification de $A$}
\begin{DispWithArrows}
  A &= \not{D_3} \left[ (D_2 \xor D_1) + \not{( \not{D_1} + D_0)} \right]
  \Cuz{T13' DeMorgan}\\
  A &= \not{D_3} \left[ (D_2 \xor D_1) + ( \not{D_1'} \not{D_0}) \right]
  \Cuz{T4 Involution}\\
  A &= \not{D_3} \left[ (D_2 \xor D_1) + ( D_1 \not{D_0}) \right]
  \Cuz{Expansion du \textit{XOR}}\\
  A &= \not{D_3} \left[ (\not{D_1}D_2 + D_1\not{D_2}) +  (D_1 \not{D_0}) \right]
  \Cuz{T7 Associativité}\\
  A &= \not{D_3} \left[ \not{D_1}D_2 + D_1\not{D_2} +  D_1 \not{D_0} \right]
  \Cuz{T8 Distributivité}\\
  A &= \not{D_3}\not{D_1}D_2 + \not{D_3}D_1\not{D_2} +  \not{D_3}\not{D_0}D_1
\end{DispWithArrows}

\subsection{Simplification de $B$}

\begin{DispWithArrows}
  B &= D_3(\not{D_1}\not{(D_2\not{D_1})})
  \Cuz{T13 DeMorgan}\\
  B &= D_3(\not{D_1}(\not{D_2}+\not{D_1'}))
  \Cuz{T4 Involution}\\
  B &= D_3(\not{D_1}(\not{D_2}+D_1))
  \Cuz{T8 Distributivité}\\
  B &= D_3(\not{D_1}\not{D_2}+\not{D_1}D_1)
  \Cuz{T5' Compléments}\\
  B &= D_3(\not{D_1}\not{D_2}+0)
  \Cuz{T1 Identité}\\
  B &= D_3(\not{D_1}\not{D_2})
  \Cuz{T7' Associativité}\\
  B &= \not{D_1}\not{D_2}D_3
\end{DispWithArrows}

\subsection{Simplification de $C$}
\begin{DispWithArrows}
  C &= (D_1\xor D_3)(D_1\xor D_2)(D_1(D_2\not{D_1} + \not{D_2}\not{D_1}))  \Cuz{T8 Distributivité}\\
    &= (D_1\xor D_3)(D_1\xor D_2)(D_1D_2\not{D_1} + D_1\not{D_2}\not{D_1}) \Cuz{T5' Compléments}\\
    &= (D_1\xor D_3)(D_1\xor D_2)(0 D_2 + 0\not{D_2}) \Cuz{T2' Élément nul}\\
    &= (D_1\xor D_3)(D_1\xor D_2)(0+0) \Cuz{T1 Identité} \\
    &= (D_1\xor D_3)(D_1\xor D_2)0 \Cuz{T2' Élément nul} \\
    &= 0
\end{DispWithArrows}

\subsection{Simplification finale}
\begin{DispWithArrows}
  R &= A+B+C \\
  R &= \not{D_3}\not{D_1}D_2 + \not{D_3}D_1\not{D_2} +  \not{D_3}\not{D_0}D_1+
       \not{D_1}\not{D_2}D_3+
       0
      \Cuz{T1 Identité}\\
  R &= \not{D_3}\not{D_1}D_2 + \not{D_3}D_1\not{D_2} + \not{D_3}\not{D_0}D_1 + \not{D_1}\not{D_2}D_3
      \Cuz{T6 Commutativité}\\
  R &=  \not{D_3}D_1\not{D_2} + \not{D_3}\not{D_0}D_1 + \not{D_1}\not{D_2}D_3 + \not{D_3}\not{D_1}D_2
      \Cuz{T6' Commutativité}\\
  R &=  \not{D_3}\not{D_2}D_1 + \not{D_3}D_1\not{D_0} + D_3\not{D_2}\not{D_1} + \not{D_3}D_2\not{D_1}
\end{DispWithArrows}

\begin{equation}
  \not{D_3}\not{D_2}D_1 + \not{D_3}D_1\not{D_0} + D_3\not{D_2}\not{D_1} + \not{D_3}D_2\not{D_1}
  =  \not{D_3}\not{D_2}D_1 + \not{D_3}D_1\not{D_0} + D_3\not{D_2}\not{D_1} + \not{D_3}D_2\not{D_1}
\end{equation}

\textbf{\textit{Q.E.D.}}

\end{document}
